% Vérifié : billet situé à 900 m de Barlette ; renvois aux ff. 41 r° et 46 v° ajoutés ; « 2e division » de Bertrancourt et destination des blessés (Varennes, < 3 km) présentées comme probables ; « usine de tissus » (lecture incertaine) nuancée ; étapes du 7th Bedfordshire « à peu de chose près » ; concordance météo et oisiveté atténuées ; date du 7 janv. 1915 précisée ; OH p. 462-463 ; passage R.A.T. abrégé (déjà dans l'introduction générale) ; phrase finale non sourcée supprimée ; guillemets droits remplacés.
% chapitre_IV.tex -- introduction du chapitre IV (fragment LaTeX, sans préambule).
% Clichés 067-104 (093, doublon de 091, ignoré) ; ff. 34 r° - 52 r°. Transcription : lecture/p2 (098-099), sinon lecture/p1.
% Sources : commentaires historiques 066-094 ; dossiers R1, R2, R3, R4, R7, R8 (hc_work/recherche).

Revenu de permission le 6~octobre 1916, Hippolyte Combes retrouve la 152nd Company au repos, loin du front, au hameau de Barlette, près de Ribeaucourt ; neuf jours plus tard, il la suit vers l'Ancre jusqu'aux briqueteries d'Albert, où elle campe pendant les dernières semaines de la bataille de la Somme. Le 2~novembre, un ordre du commandant du train le détache avec le lieutenant Curtis à Acheux, pour acheter de la paille et du foin. Le chapitre fait donc passer l'interprète d'une vie réglée sur les étapes et les corvées de la compagnie à un service de bureau et de tournées en automobile (\textit{Feuillets de Route}, chap.~\textsc{x} : «~Randonnées quotidiennes en auto~», annexe, p.~\pageref{annexe:doc6}), à l'arrière d'une bataille qu'il entend sans la voir, et qui s'achève le 18~novembre.

\subsection*{Le carnet}

À neuf cents mètres de Barlette, hameau de quelques maisons, Hippolyte Combes loge chez « Mme Ve — », dont il laisse le nom en blanc : une « chambre très belle \& spacieuse » (f.~34~r\textsuperscript{o}), ancien salon d'une « vieille gd-mère, extt sympathique » (f.~34~v\textsuperscript{o}). Il fait viser son titre de permission, puis touche sa solde au bureau de la Mission française de Bernaville (ff.~34~r\textsuperscript{o} et 35~r\textsuperscript{o}). Le service se réduit au fourrage vert : il conduit les fourragères dans les champs de luzerne, « assis s/ 1 botte de luzerne » pendant le chargement (f.~36~r\textsuperscript{o}), ou, sous une pluie battante, mène des voitures dans un champ de luzerne derrière le grand bois de Ribeaucourt (ff.~36~v\textsuperscript{o}--37~r\textsuperscript{o}). Ses loisirs vont à la forêt : « Resté lgs ds cet océan de verdure » (f.~35~v\textsuperscript{o}). Le 13~octobre au dîner, il apprend le départ : « Regret de partir ! » (f.~37~r\textsuperscript{o}).

Le mouvement dure cinq jours. Le 15, la colonne gagne Beauval, « remplie de troupes » (f.~37~v\textsuperscript{o}) ; le 16, par Beauquesne et Puchevillers, Warloy-Baillon, « vraie fourmilière de tommies », où l'on campe sous la tente (f.~38~r\textsuperscript{o}) ; le 17, un champ dévasté entre Bouzincourt et Hédauville, où Hippolyte Combes dort sur un fauteuil de toile dans la tente du mess, sous la pluie : « Aucun confort » (f.~39~r\textsuperscript{o}). Le 18, un camion chargé d'obus le mène à Albert, où il s'arrête devant la « basilique ravagée par les obus » (f.~39~v\textsuperscript{o}). Sur la route, il voit passer nuit et jour camions et caissons, « Autant de mouvement que le long d'1 gde artère de capitale » (f.~40~r\textsuperscript{o}). Le 19, la compagnie s'installe dans la boue du \textit{brickfield area} ; Hippolyte Combes trouve une chambre dans une rue voisine (f.~40~v\textsuperscript{o}).

À Albert, du 19~octobre au 2~novembre, le carnet note d'abord le canon : « Duel d'artillerie formidable » le 20, vitres et portes qui tremblent, éclairs des pièces à l'horizon la nuit (f.~41~r\textsuperscript{o}). Le 22, Hippolyte Combes décrit la Vierge de la basilique, qui, « horiz\textsuperscript{le}, penche Jésus bras ouverts au-dessus place », et l'adjudant Schumacher lui apprend que « les 2/3 des interprètes, jeunes R.A.T. seront renvoyés regiments frs » (f.~42~r\textsuperscript{o}) ; le 23, il visite les abris allemands de La Boisselle et voit l'« Immense cratère s/ la droite » : « Impression funèbre et désolation » (f.~42~v\textsuperscript{o}). Le 24, l'adjudant [Lépin~?] et le lieutenant Pizzi l'assurent qu'il ne quittera pas la Mission (f.~43~r\textsuperscript{o}). Suivent des jours de pluie et d'inaction : « Vie d'1 extrême monotonie » (f.~43~r\textsuperscript{o}), « Aucune occupation comme interprète » (f.~43~v\textsuperscript{o}). Il veille avec Abel, le plus jeune fils de ses logeurs (f.~41~v\textsuperscript{o}), et écoute l'aîné, maréchal des logis dans l'artillerie belge, raconter ses combats (f.~44~r\textsuperscript{o}). Le 29, une course en auto à Acheux et Senlis, « enquête au sujet de straw \& hay disponibles » (\textit{ibid.}), annonce le changement. Le jour de la Toussaint, après une visite au cimetière ravagé et une prise d'armes, où un général écossais décore des soldats distingués dans l'offensive de la Somme, Curtis lui apporte l'ordre : « Ns sommes détachés de la Coy Transport » (f.~45~r\textsuperscript{o}).

Sur le conseil du bureau de la Mission, Hippolyte Combes demande un ordre écrit ; au \textit{Train Office}, « Major Hughes tempêté », puis le colonel Grose signe (f.~45~v\textsuperscript{o}). Le soir même, à Raincheval, il voit le « major Levy, of the D.O.S.T. » (f.~46~r\textsuperscript{o}). Délogé le 5 de son premier billet par le \textit{Town Major}, il s'installe au n\textsuperscript{o}~1 de la route de Léalvillers, chez un propriétaire « très sympathiq\textsuperscript{ue} », et il est convoqué le soir même par le capitaine Huguet au « bureau mission du corps » (f.~47~r\textsuperscript{o}). Le premier achat de paille et de foin que note le carnet a lieu le 6, à Coisy ; le même jour, il « Regrette d'avoir quitté 152 Coy » et « Cherche moyen de ne pas quitter 18\textsuperscript{ie} divon \& ma Coy » (f.~47~v\textsuperscript{o}). Dès lors, le carnet aligne les courses en auto, de Forceville à Pierregot et jusqu'à Doullens et Frévent (ff.~48~r\textsuperscript{o}--52~r\textsuperscript{o}), et les passages au bureau de la Mission d'Albert, qui le paie encore ; le 9, il prend le thé au mess de son ancienne compagnie (f.~48~r\textsuperscript{o}). Le 10, seul au bureau, il lit et écrit près du feu : « Impression toute nouvelle » (f.~48~v\textsuperscript{o}). Le 11, à Varennes, il rencontre l'interprète d'une batterie, professeur au lycée de Toulon ; le soir, Acheux est bombardé et l'on se réfugie à la cave (\textit{ibid.}).

La bataille de l'Ancre ne paraît dans le carnet que par le bruit et par ses suites : « très fort bombardt » du côté de Beaumont et passage de blessés sur des fourgons, le 14 (f.~49~r\textsuperscript{o}) ; « Vitres et fenêtres ne cessent pas de trembler », le 16 (f.~49~v\textsuperscript{o}) ; canonnade de l'aube du 18, où « on dirait énormes vagues déferlant s/ rochers », sous la « 1\textsuperscript{ère} neige » (f.~50~r\textsuperscript{o}). Le 20, de sa fenêtre, il regarde les troupes qui reviennent des tranchées et les « cortèges de prisonniers allemands » (f.~50~v\textsuperscript{o}). Le reste du mois est fait de boue, de froid et de petites misères : une « remarque +tot desobligeante de Curtis » le 21 (\textit{ibid.}), une indisposition le 22 (f.~51~r\textsuperscript{o}), plus de vingt-huit obus sur Acheux le 25, « Plusrs tués et blessés », et de longues heures au bureau, assis sur une caisse de bidons de pétrole : « Triste vie » (f.~51~v\textsuperscript{o}). Restent les rencontres : le lieutenant Walsh et son interprète, directeur d'une usine (de tissus~?) près de Paris, à Hérissart (ff.~49~v\textsuperscript{o} et 51~v\textsuperscript{o}), un thé au château de Martinsart, avec des « Gros canons, tt proche, en action » (f.~52~r\textsuperscript{o}), une conférence sur la Russie dans la baraque de la Y.M.C.A. (f.~50~r\textsuperscript{o}). Les comptes restent modestes : cartes postales presque chaque jour, 29~F puis 47,70~F pour le mess de septembre et d'octobre (ff.~34~v\textsuperscript{o} et 46~r\textsuperscript{o}), un casque peint en kaki, un déjeuner de 7~F à l'hôtel d'Amiens (f.~49~r\textsuperscript{o}). Le 30, Hippolyte Combes dîne pour la première fois au nouveau mess, dans la chambre de Curtis (f.~52~r\textsuperscript{o}).

\subsection*{Le contexte}

Les déplacements d'Hippolyte Combes suivent ceux de la 54th Brigade, que servait la 152nd Company. Relevée après Thiepval et la Schwaben Redoubt, la 18th Division se repose depuis le 2~octobre « in the Candas area »\footnote{G.~H.~F.~Nichols, \textit{The 18th Division in the Great War}, Édimbourg et Londres, Blackwood, 1922, p.~126.}. Le 7th Bedfordshire cantonne à Vacquerie et à Gorges, à quelques kilomètres de Barlette ; il gagne Beauval le 15, le bois de Vadencourt le 16, Bouzincourt le 17 et Albert le 19, soit, à peu de chose près, les étapes du carnet\footnote{Journal de marche (\textit{War Diary}) du 7th (Service) Battalion, Bedfordshire Regiment, octobre 1916, Bedfordshire Archives, X550/8/1/1610 (transcription en ligne). Le tableau de Nichols (\textit{op.~cit.}, p.~108--110), qui attache la 152nd Company à la 55th Brigade, est contredit par les ordres de marche du même bataillon (X550/8/1/1609 et 1704).}. La division revient pour l'attaque de Regina Trench, lancée par le II Corps le 21~octobre à 12~h~06, où sa 53rd Brigade prend sa part de la tranchée : Hippolyte Combes en entend la préparation le 20 (f.~41~r\textsuperscript{o})\footnote{W.~Miles, \textit{Military Operations, France and Belgium, 1916}, t.~II, Londres, Macmillan, 1938 (ci-après \textit{OH 1916}, II), p.~462--463 ; Nichols, \textit{op.~cit.}, p.~131.}. Les notes météorologiques du carnet concordent avec l'histoire officielle : l'orage du 19 fait remettre l'attaque de vingt-quatre heures, et la pluie revenue le 24 fait reporter de date en date l'opération principale sur l'Ancre\footnote{\textit{OH 1916}, II, p.~461--462.}. Les troupes l'attendent pendant des semaines, « froides et souvent mouillées jusqu'aux os », sous la menace d'une attaque « constamment remise mais jamais annulée »\footnote{\textit{Ibid.}, p.~480. Nichols (\textit{op.~cit.}, p.~130) résume : « October and November 1916 were months in which the Staff proposed and the weather disposed ».}. Le 30~octobre, la Reserve Army du général Gough prend le nom de Fifth Army\footnote{A.~F.~Becke, \textit{Order of Battle of Divisions}, part~3A, \textit{New Army Divisions (9--26)}, Londres, HMSO, 1938, p.~84.}.

Albert, qui comptait 6\,742 habitants avant la guerre\footnote{J.~E.~Edmonds, \textit{Military Operations, France and Belgium, 1916}, t.~I, Londres, Macmillan, 1932 (ci-après \textit{OH 1916}, I), p.~247.}, n'est plus qu'une ville de passage aux maisons vides, ce qui explique sans doute l'oisiveté de l'interprète, d'abord « trait d'union entre les Anglais et les habitants du pays »\footnote{Préface de F.~de Tessan à Fernand-Laurent, \textit{Chez nos alliés britanniques. Notes et souvenirs d'un interprète}, Paris, Boivin, 1917, p.~XIV.}. Les obus qu'Hippolyte Combes note le 24 et le 31~octobre, puis le 9~novembre, nuancent le guide Michelin, selon lequel la ville était « enfin hors de portée » des canons allemands en octobre 1916\footnote{Michelin, \textit{Illustrated Michelin Guides to the Battlefields (1914--1918). The Somme}, vol.~1, \textit{The First Battle of the Somme (1916--1917)}, Clermont-Ferrand, Michelin, 1919, p.~31--35 ; ff.~43~r\textsuperscript{o}, 44~v\textsuperscript{o} et 48~r\textsuperscript{o}.}. Hippolyte Combes date de 1914 les obus qui ont atteint la tour de la basilique ; c'est le 7~janvier 1915 que le dôme fut crevé, puis, le 15, que la coupole vola en éclats et que la statue s'inclina, avant de se coucher à l'horizontale le soir même\footnote{Diocèse d'Amiens, « Notre-Dame de Brebières ou la “Vierge penchée” », 13~mars 2018, \url{https://amiens.catholique.fr} ; base Mérimée, notice PA80000002.}.

Acheux, où il arrive le 2~novembre, est depuis le 2~octobre le siège du quartier général du V Corps, chargé de l'effort principal au nord de l'Ancre\footnote{C.~Baker, \textit{The Long, Long Trail}, « The British V Corps in France and Flanders », copie de la Wayback Machine du 5~août 2020, \url{https://web.archive.org/web/20200805133825/https://www.longlongtrail.co.uk/army/other-aspects-of-order-of-battle/the-british-corps-headquarters-in-france-and-flanders/the-british-v-corps-in-france-and-flanders/}.} ; le « bureau mission du corps » du 5~novembre est vraisemblablement le sien. À Bertrancourt, le 4, Hippolyte Combes est passé au « bureau mission fr. 2\textsuperscript{e} division » (f.~46~v\textsuperscript{o}), sans doute celui de la 2nd Division, l'une des quatre divisions d'assaut du corps\footnote{\textit{OH 1916}, II, p.~477.}, où il retrouve l'interprète Cordonnier « que connu à Bethune » lors de son service de mai--juillet 1915 auprès des \textit{1\textsuperscript{st} Irish Guards}, alors rattachés à cette même division (\textit{Feuillets de Route}, chap.~\textsc{v}--\textsc{vi}, p.~\pageref{annexe:doc6}). L'assaut est donné le 13 : Beaumont-Hamel tombe le jour même, Beaucourt le 14, et la bataille rapporte plus de 7\,200 prisonniers\footnote{J.~H.~Boraston (éd.), \textit{Sir Douglas Haig's Despatches (December 1915 -- April 1919)}, Londres et Toronto, Dent, 1919, p.~49--51.}. Le 18, l'assaut final part à 6~h~10 « dans un tourbillon de grésil qui se changea ensuite en pluie », après la première neige de l'hiver ; la 55th Brigade de la 18th Division y attaque Desire Trench, et cette date marque officiellement la fin des batailles de la Somme\footnote{\textit{OH 1916}, II, p.~xliv et 514--517.}. Une partie des blessés qu'Hippolyte Combes voit passer va sans doute aux hôpitaux d'évacuation de Varennes, à moins de trois kilomètres au sud d'Acheux, où les 4th et 11th Casualty Clearing Stations sont installées depuis octobre\footnote{Commonwealth War Graves Commission, notice « Varennes Military Cemetery », \url{https://www.cwgc.org} ; l'histoire officielle en compte trois dès novembre (\textit{OH 1916}, II, p.~541).}. Les bombardements d'Acheux du 11 et du 25~novembre ne sont pas mentionnés dans les sources consultées.

Le service auquel Hippolyte Combes est affecté répond à une contrainte nouvelle. En novembre 1916, la crise du tonnage conduit le \textit{War Office} à proposer d'« étendre l'achat local » du fourrage\footnote{\textit{OH 1916}, II, p.~550.} ; la paille sert à la fois de complément de ration pour les chevaux et de couchage, à raison de quatre livres par homme et par semaine en hiver, et chaque armée l'achète par son \textit{purchase board}\footnote{F.~K.~Puckle, \textit{The Army Service Corps of the British Army and the Organization of the Transport and Transportation at the Front in France}, Washington, Government Printing Office, 1917, p.~8 et 28.}. L'ordre de détachement est signé par le lieutenant-colonel Daniel Charles Evans Grose, commandant du train de la division\footnote{Nichols, \textit{op.~cit.}, p.~107 ; \textit{Supplement to the London Gazette}, n\textsuperscript{o}~30111, 4~juin 1917, p.~5470.}. Le major Levy est peut-être Walter Henry Levy, major temporaire de l'Army Service Corps depuis le 1\textsuperscript{er}~mai 1916, et le major Hughes, Herbert Francis Hughes, promu dans la même liste ; ce ne sont que des hypothèses, et Curtis, faute de prénom, reste non identifié\footnote{\textit{Supplement to the London Gazette}, n\textsuperscript{o}~29585, 17~mai 1916, p.~4941 ; The National Archives (Kew), catalogue Discovery, dossiers d'officiers WO~339 et WO~374, où figurent onze officiers de l'ASC nommés Curtis.}. La rumeur du 22~octobre touche Hippolyte Combes de près, puisqu'il appartient lui-même, selon toute vraisemblance, à la réserve de l'armée territoriale depuis le 1\textsuperscript{er}~novembre 1916 (voir l'introduction)\footnote{Fiche matricule d'Hippolyte Combes, recrutement de Toulouse, classe 1896, n\textsuperscript{o}~911, archives familiales.}. Aucun texte relatif à cette réduction des interprètes n'a été retrouvé\footnote{Les dossiers de la Mission militaire française attachée à l'armée britannique (Service historique de la Défense, GR~16~N~509 et 17~N~447--448) n'ont pas été consultés.}.

La fin de novembre prépare le chapitre suivant. Le 29~novembre, le V Corps quitte Acheux pour Doullens et le XIII Corps prend sa place\footnote{Baker, notice citée, et « The British XIII Corps in France and Flanders », \url{https://www.longlongtrail.co.uk} (consultée le 18~octobre 2026).} ; la 18th Division, relevée, part se reposer dans la région d'Abbeville\footnote{Nichols, \textit{op.~cit.}, p.~136 et 139.}. Hippolyte Combes, lui, reste à Acheux. Le 1\textsuperscript{er}~décembre, il est « pris en charge » par le bureau de la Mission du XIII Corps et note aussitôt : « Perspective réintégration prochaine ds régiment d'origine ! » ; le 2, le capitaine Williams arrive pour remplacer Curtis (f.~52~v\textsuperscript{o}).
